\section{Passing Through the Wrappings: Modern Hypotheses and the Boundary of the Text}
\label{sec:modern-hypotheses}

A single sentence in the Catechism of the Council of Trent opens the door to a notable modern discussion of John 20:6--7. Teaching on the Third Article of the Creed, the Roman Catechism sets three mysteries of Christ's body beside one another:
\begin{quote}
``He is born of his Mother without any diminution of her maternal virginity; and as he afterwards went forth from the sepulchre whilst it was closed and sealed, and entered the room in which his disciples were assembled, `the doors being shut' (John xx. 19); or, not to depart from natural events which we witness every day, as the rays of the sun penetrate, without breaking or injuring, in the least, the substance of glass; after a like, but more incomprehensible manner, did Jesus Christ come forth from his mother's womb without injury to her maternal virginity.''\endnote{Catechism of the Council of Trent, Part I, Article III, trans. J. Donovan (Maynooth, 1829; Baltimore: Lucas Brothers), p. 41. Compare St. John Chrysostom as excerpted in St. Thomas Aquinas, \emph{Catena Aurea in Ioannem}, cap. 20 (Oxford, 1874), p. 593: ``Our Lord rose while the stone and seal were still on the sepulchre. But as it was necessary that others should be certified of this, the sepulchre is opened after the resurrection, and so the fact confirmed.''}
\end{quote}

If the risen body of the Lord passed out of the sealed tomb before the angel rolled back the stone---and later stood among the Apostles while the doors remained shut---did He need to unwind the myrrh-stiffened burial bandages by hand at all, or did His glorified body pass straight through the linen cloths themselves?

\subsection{Henry Latham's Hypothesis: Collapsed Wrappings and an Annular Head-Cloth}

In 1901, Henry Latham, Master of Trinity Hall, Cambridge, published \emph{The Risen Master}, devoting two meticulous chapters to the geometry of the cloths in John 20:5--8. Unlike the later chain-letter story, Latham's proposal grew entirely out of close attention to Greek grammar and Palestinian tomb architecture.

Latham pictured a rock-cut tomb with a burial bench along one side, having a slightly raised stone ledge at the head end. On the lower portion of the bench lay the \latin{othonia}, weighted down with Nicodemus's myrrh and aloes. When the risen body passed out of them, Latham argued, the wrappings did not get uncoiled or scattered; they simply subsided flat onto the slab under the weight of the spices (\emph{The Risen Master}, 1905 reprint, p. 34):
\begin{quote}
``In this recess, on the lower part of the ledge, lay the grave-clothes. They were in no disorder, they were just as they were when Joseph and others had wrapped them round the body of the Lord, only they were lying flat, fold over fold, for the body was gone. On the raised part of the ledge at the far end, all by itself, was the napkin that had gone round the head; this was not lying flat, but was standing up a little, retaining the twirled form which had been given it when it had been twined round the head of the Lord. Nothing in the place gave any sign of the touch of human hands: the body had been embedded in the powdered aloes and myrrh, but of this there was not a trace; the spice remained enclosed by the cloths between which it had been placed when the body was laid on the slab.''
\end{quote}

Latham rested his interpretation of the head-cloth directly on the New Testament usage of \latin{entyliss\={o}} (p. 43):
\begin{quote}
``Now the word used for rolled up is employed (Matthew xxvii. 59, and Luke xxiii. 53) to express the wrapping of the Lord's body in the linen cloths, and it implies here that the napkin had been wrapped round the head and partially retained the annular form thus given it; I take it to mean that it was not folded so as to lie flat like the cloths. Now the Greek says that the napkin was not lying with the linen cloths, but rolled up apart in a place by itself. Our Revised Version leaves out the word apart, but what the Evangelist means I believe to be, that the napkin was separated from the rest of the clothes, and lay by itself on the raised part of the ledge, where there was nothing else.''
\end{quote}

Latham's reconstruction has an undeniable dramatic power. If the body-cloths lay collapsed in their wrapped coils with the spices still inside, and the head-cloth retained its ring-like (\latin{annular}) shape a short distance away on the headrest, then the beloved disciple needed no detective deduction about how long a thief would spend untying bandages: a single glance inside the chamber showed that the body had passed straight through the grave-clothes without disturbing them.

\subsection{Brooke Foss Westcott's Nuanced Verdict}

Bishop Brooke Foss Westcott, whose monumental commentary on St. John remained a standard work of Greek scholarship, read Latham's proposal with deep interest---and drew a sharp grammatical distinction that Latham had missed. In the posthumous 1908 two-volume Greek-text edition prepared by his son Arthur Westcott, the editor reproduces two of his father's later notes on John 20:7--8 (vol. 2, p. 339). In a private letter to the Rev. H. O. Barratt, Bishop Westcott wrote:
\begin{quote}
``The explanation which you give of John xx. 8 (passed through the heavy wrappings as He later passed through doors) is, I think, substantially correct. \dots\ I did not agree with his (sc. Mr. Latham's) interpretation of the participle \dots\ Perhaps the separate treatment of the face cloth suggests the action of the living Lord.''
\end{quote}
In a second marginal note printed immediately below it, Westcott distilled his own hybrid interpretation into two terse sentences:
\begin{quote}
``The undisturbed grave-cloths show that the Lord had risen through and out of them. The face cloth carefully rolled up, the action of the living Lord.''\endnote{Brooke Foss Westcott, \emph{The Gospel according to St. John: The Greek Text with Introduction and Notes}, ed. Arthur Westcott, vol. 2 (London: John Murray, 1908), p. 339. These remarks stand in Arthur Westcott's signed editorial footnotes (\latin{A. W.}), recording Bishop Westcott's private letter and later marginal notes alongside his published commentary.}
\end{quote}

Westcott's reservation goes straight to the heart of the Greek sentence. Verse 7 builds a direct antithesis: \latin{ou meta t\={o}n othoni\={o}n keimenon, alla ch\={o}ris entetyligmenon eis hena topon}---``not lying with the linen cloths, but apart, rolled up into one place.'' If the \latin{soudarion} simply remained sitting untouched where Nicodemus had wrapped it around the head, only inches from the neck of the body-cloths, it is hard to see why St. John would emphasize so strongly that it was \latin{ou meta t\={o}n othoni\={o}n} (``not with the linen cloths'') and \latin{ch\={o}ris \dots\ eis hena topon} (``apart \dots\ into one place''). The preposition \latin{eis} (``into'') with \latin{ch\={o}ris} (``apart'') fits naturally with a cloth set deliberately into a place of its own---whether rolled up by the living Lord, as Chrysostom and Westcott supposed, or set aside by the angels attending the tomb.

\subsection{Three Reconstructions and the Honest Limit of the Words}

Placed side by side, the three serious interpretations grounded in the Greek text reveal exactly where the text speaks firmly and where disciplined inference begins:

\begin{itemize}
\item \textbf{Unhurried removal (St. John Chrysostom, St. Thomas Aquinas, Cornelius a Lapide):} The linen cloths lie in one spot and the head-cloth lies rolled up apart in another, demonstrating calm, deliberate order incompatible with grave-robbers working in fear and haste over a myrrh-glued corpse.
\item \textbf{Untouched collapsed wrappings (Henry Latham):} The glorified body passed out of the spice-weighted \latin{othonia}, leaving them collapsed flat on the bench and the \latin{soudarion} retaining its coiled ring-shape on the separate headrest.
\item \textbf{Hybrid (Brooke Foss Westcott):} The body passed directly out of the undisturbed \latin{othonia}, while the separate, rolled-up \latin{soudarion} reflects the deliberate action of the living Lord.
\end{itemize}

Every one of these three readings respects the Greek grammar of John 20:5--8, and none invents a phantom first-century custom. Yet the Evangelist himself leaves the exact mechanics unstated. St. John does not say whether the cloths were collapsed flat or unwound; he does not say whose hands rolled up the \latin{soudarion}; he does not even state explicitly whether the burial bench had a raised headrest. What the Greek text establishes beyond doubt is the visible result that met Simon Peter and the beloved disciple inside the rock: no frantic struggle, no torn wrappings strewn across the threshold, no stripped corpse dragged into the dark, but ordered linen cloths lying in place and a head-cloth rolled up apart in one spot.

\subsection{Material Relics and the Boundary of Exegesis}

Catholic devotion and modern historical inquiry also encounter two famous linen relics associated with the Passion and Burial of Christ: the Shroud of Turin (\latin{sind\={o}n}) and the Sudarium of Oviedo (\latin{sudarium}). Both textiles have generated a vast, specialized literature involving textile weaving, bloodstain geometry, pollen analysis, medieval provenance, and radiocarbon controversy. That entire field belongs to forensic archaeology, art history, and material conservation rather than the philological and patristic exegesis of John 20:3--10 and Luke 24:12. The words of the Gospels stand on their own historical and theological foundation, independent of any later textile investigation, just as the faith of the Church in the bodily Resurrection of Christ rests on the apostolic testimony to the risen Lord rather than on any surviving linen cloth.
